\section{Presentación y objetivos}
\label{sec:presentacion}

\subsection{Presentación}
\label{sec:presentacion-general}

Este informe reúne y describe los datos disponibles de la estación de medida
\EstacionInforme{} correspondientes al periodo analizado. Su finalidad es
ofrecer una visión ordenada de la disponibilidad de observaciones y de las
características de las propiedades atmosféricas registradas.

\subsection{Objetivos}
\label{sec:objetivos}

Los objetivos del informe son:

\begin{enumerate}
  \item Documentar los instrumentos y los productos utilizados en el análisis.
  \item Identificar los periodos con datos disponibles y cuantificar su cobertura.
  \item Describir la distribución mensual de las variables de precipitación y de
        las propiedades estimadas de los hidrometeoros que permitan los datos
        disponibles.
  \item Proporcionar figuras y tablas reproducibles que faciliten la revisión y
        la comparación de los meses observados.
\end{enumerate}

\subsection{Perspectiva científica y alcance}
\label{sec:alcance}

El estudio tiene un carácter descriptivo y exploratorio. Las observaciones de
superficie y las medidas del radar aportan perspectivas complementarias sobre
la precipitación, aunque sus volúmenes de medida, métodos de estimación y
limitaciones son diferentes.

Las diferencias entre meses se describirán teniendo en cuenta la cobertura y
los criterios de calidad.

El informe no pretende establecer una climatología de la estación ni atribuir
causas a las variaciones observadas. Su propósito es caracterizar los datos
efectivamente disponibles en el periodo analizado y dejar explícitas las
condiciones bajo las que pueden interpretarse sus resultados.
